% Proof fragment for the existing flat-coarse-rank dossier; no standalone paper.
% All norms below use counting measure. Analytic proof, not Lean verified.
\section{A finite proof of the critical rectangle estimate}
\label{sec:self-contained-box-sobolev}
This supplement derives the balanced-cube endpoint used by the rectangle
adapter, rather than importing a local Sobolev estimate as an external input.
The method is classical: finite entropy proves the projection inequality,
which yields relative edge isoperimetry, layer sets, and a power argument.
We claim neither a new general Sobolev principle nor a solution of the open
critical rectangular-partition problem. This is part of the existing dossier.

Fix $D\geq3$, and put
\[
 \begin{gathered}
 p=D/2,\quad q=\frac{2D}{D+2},\quad q'=\frac{2D}{D-2},\\
 \alpha=\frac{2(D-1)}D,\quad c_D=D(1-2^{-1/D}),\quad
 K_D=\frac{2\alpha D^{1/q'}}{c_D}.
 \end{gathered}
\]
An internal axis edge is an unordered adjacent pair, counted once. For its
positive-axis orientation $e=(x,x+e_i)$, write $du(e)=u(x+e_i)-u(x)$.

\subsection{Finite projection inequality, with proof}
For a nonempty finite $A\subset\prod_{j=1}^D X_j$, let $\pi_i$ delete
coordinate $i$. Then
\begin{equation}\label{eq:finite-projection}
 |A|^{D-1}\leq\prod_{i=1}^D|\pi_i A|.
\end{equation}
Here is a finite proof, including the entropy facts needed for this use.
For a finite probability vector $r$, define $H(r)=-\sum r\log r$ with
$0\log0=0$. Concavity follows from $(-t\log t)''=-1/t<0$ and continuity at
zero. Consequently mixing conditional distributions cannot decrease their
average entropy: $H(Y\mid Z)\geq H(Y\mid Z,W)$. Factoring finite joint
probabilities proves the chain rule
$H(X)=\sum_jH(X_j\mid X_1,\ldots,X_{j-1})$.
Also $H(Y)\leq\log|\operatorname{supp}Y|$: Jensen's inequality for $\log$
gives $\sum_y r_y\log(1/r_y)\leq\log\sum_y r_y/r_y$ on the positive support.
For $X$ uniform on $A$, $H(X)=\log|A|$, while
\[
 \begin{split}
 \sum_i H(X_{-i})
 &=\sum_i\sum_{j\ne i}H(X_j\mid X_{<j}\setminus X_i)\\
 &\geq\sum_i\sum_{j\ne i}H(X_j\mid X_{<j})
 =(D-1)H(X).
 \end{split}
\]
The notation in the conditioning set means omit $X_i$ only when $i<j$.
Bounding each $H(X_{-i})$ by $\log|\pi_i A|$ and exponentiating proves
\eqref{eq:finite-projection}. No projection inequality is an unproved
hypothesis in the sequel.

\subsection{Relative isoperimetry on a finite cube}
Let $Q_s=\{0,\ldots,s-1\}^D$, $s\geq2$, and let $\partial_E A$ consist
of internal cube edges crossing $A$ and its complement. If
$0<|A|\leq |Q_s|/2$, then
\begin{equation}\label{eq:relative-edge-iso}
 |\partial_E A|\geq c_D|A|^{(D-1)/D}.
\end{equation}
Indeed, put $a=|A|$, $P_i=|\pi_i A|$. At most $a/s$ occupied fibers in
direction $i$ are full. Every occupied nonfull fiber contains an internal
crossing edge, because the fiber is a connected interval. Edges in distinct
fibers/directions are different. Thus projection inequality and AM--GM give
\[
 |\partial_E A|\geq\sum_i(P_i-a/s)
 \geq D a^{(D-1)/D}-Da/s
 \geq D(1-2^{-1/D})a^{(D-1)/D}.
\]
Empty $A$ satisfies the same inequality trivially. The boundary is the EDGE
boundary, not the vertex boundary in a separate graph theorem.

\subsection{Endpoint counting-norm Sobolev--Poincare inequality}
For every real $u$ on $Q_s$, including $s=1$,
\begin{equation}\label{eq:self-contained-cube}
 \sum_{Q_s}|u-\bar u_{Q_s}|^2
 \leq K_D^2\left(\sum_{e\subset Q_s}|du(e)|^q\right)^{2/q}.
\end{equation}
For $s=1$ both sides vanish. For a nonnegative $v$ supported on at most
half the cube, layer-set decomposition and the triangle inequality in
$\ell^{D/(D-1)}$, followed by \eqref{eq:relative-edge-iso}, imply
\[
 \|v\|_{D/(D-1)}
 \leq\int_0^\infty|\{v>t\}|^{(D-1)/D}\,dt
 \leq c_D^{-1}\int_0^\infty|\partial_E\{v>t\}|\,dt
 =c_D^{-1}\sum_e|dv(e)|.
\]
The integrals are finite sums of step lengths between values of $v$;
the last equality follows edgewise from the length of the interval
between its two endpoint values. Apply this inequality to $v^\alpha$.
For $a,b\geq0$, the fundamental theorem of calculus gives
\[
 |a^\alpha-b^\alpha|
 \leq\alpha|a-b|(a^{\alpha-1}+b^{\alpha-1}).
\]
Since $\alpha D/(D-1)=2$ and $(\alpha-1)q'=2$, Holder and the degree
bound $\deg_{Q_s}(x)\leq2D$ give
\[
 \begin{split}
 \|v\|_2^\alpha
 &\leq\frac\alpha{c_D}\Big(\sum_e|dv(e)|^q\Big)^{1/q}
 \Big(\sum_{e=\{x,y\}}(v(x)^{\alpha-1}+v(y)^{\alpha-1})^{q'}\Big)^{1/q'}\\
 &\leq\frac{2\alpha D^{1/q'}}{c_D}
 \Big(\sum_e|dv(e)|^q\Big)^{1/q}\|v\|_2^{\alpha-1}.
 \end{split}
\]
The second inequality uses $(a+b)^{q'}\leq2^{q'-1}(a^{q'}+b^{q'})$ and
$\sum_e(v(x)^2+v(y)^2)\leq2D\|v\|_2^2$.
If $v=0$ the conclusion is immediate; otherwise division proves
$\|v\|_2\leq K_D\|dv\|_q$.

Choose a counting median $m$ of $u$ and set $v_+=(u-m)_+$,
$v_-=(m-u)_+$. Both supports have size at most half the cube. For each edge,
\[
 |dv_+(e)|^q+|dv_-(e)|^q\leq|du(e)|^q.
\]
For endpoint values on the same side of $m$ this is equality in the
nonzero part; on opposite sides it is $a^q+b^q\leq(a+b)^q$.
Since $2/q\geq1$, addition of the two squared estimates gives
\[
 \|u-\bar u_{Q_s}\|_2^2\leq\|u-m\|_2^2
 \leq K_D^2\big(\|dv_+\|_q^q+\|dv_-\|_q^q\big)^{2/q}
 \leq K_D^2\|du\|_q^2.
\]
This proves \eqref{eq:self-contained-cube} without an external graph,
ball-adapter, coefficient, or boundary-condition hypothesis.

\subsection{Reflection with all aspect-ratio multiplicities}
Let $R=\prod_i\{0,\ldots,l_i-1\}$, $V=\prod_i l_i$,
$L=\max_i l_i$, and $\kappa=L^2/V^{2/D}$. If $V=1$ the desired variance
bound is zero. Otherwise $L\geq2$. Put $s=2L+1$, $r=s^D/V$ and fold
$Q_s$ coordinatewise onto $R$ by the endpoint-repeating pattern of period
$2l_i$,
\[
 0,1,\ldots,l_i-1,l_i-1,\ldots,1,0,0,1,\ldots.
\]
For $l_i=1$ the fold is constant. Since $s/(2l_i)>1$, each original
coordinate value has at least $2\lfloor s/(2l_i)\rfloor\geq s/(2l_i)$
preimages and at most $s/l_i+2\leq3s/l_i$ preimages. An original internal
edge has at most the same upper count in its own direction; each full
period traverses it twice. Adjacent folded values either coincide or are
an original internal edge, possibly with reversed orientation. Thus for
the unweighted extension $\widetilde u$,
\[
 \operatorname{Var}_{Q_s}(\widetilde u)\geq2^{-D}r\operatorname{Var}_R(u),
 \qquad
 \sum_{e\subset Q_s}|d\widetilde u(e)|^q
 \leq3^D r\sum_{e\subset R}|du(e)|^q.
\]
The variance inequality holds by applying the preimage lower bound at the
mean of $\widetilde u$, and then minimizing over constants on $R$.
Apply \eqref{eq:self-contained-cube}. Since $2/q-1=2/D$ and $s\leq3L$,
\begin{equation}\label{eq:rectangle-unweighted-derived}
 \operatorname{Var}_R(u)\leq A_D\kappa
 \left(\sum_{e\subset R}|du(e)|^q\right)^{2/q},\qquad
 A_D=K_D^2 2^D 3^{2D/q+2}.
\end{equation}
The factor $r$ is not bounded independently of aspect ratio; it was kept
until it became $s^2/V^{2/D}$. Singleton coordinate sides cause no problem.

\subsection{Positive original edge weights and field coordinates}
For arbitrary positive ORIGINAL edge coefficients $b_e$, Holder on $R$
with exponents $2/q$ and $2/(2-q)$ gives
\[
 \left(\sum_e|du|^q\right)^{2/q}
 \leq\Big(\sum_e b_e|du|^2\Big)
 \Big(\sum_e b_e^{-q/(2-q)}\Big)^{(2-q)/q}.
\]
Here $q/(2-q)=D/2=p$. Consequently
\begin{equation}\label{eq:rectangle-weighted-derived}
 \operatorname{Var}_R(u)\leq A_D\kappa
 \Big(\sum_{e\subset R}b_e^{-D/2}\Big)^{2/D}
 \sum_{e\subset R}b_e|du(e)|^2.
\end{equation}
No reflected weights were introduced: reflection was completed before
this Holder step. For $b_{(x,x+e_i)}=a(x)>0$, each vertex supplies at most
$D$ internal starting edges. Thus the inverse edge mass is at most
$D\sum_Ra^{-D/2}$, proving exactly the scalar-start estimate
\[
 \operatorname{Var}_R(u)\leq A_D D^{2/D}\kappa
 \Big(\sum_Ra(x)^{-D/2}\Big)^{2/D}\mathcal E_{a,R}(u).
\]
For any fixed finite real internal/field dimension, apply this scalar
inequality to each coordinate and sum; its coefficient factor is common.
There is no additional internal-dimension factor and no replacement of
original edges by enlarged-region energy.

\subsection{Necessity of geometry and status}
For uniform $a>0$ and a longest coordinate of length $L\geq2$, put
$u(x)=x_i-(L-1)/2$. Direct sums give
\[
 \operatorname{Var}_R(u)=V(L^2-1)/12,\qquad
 \mathcal E_{a,R}(u)=aV(L-1)/L,
\]
so their ratio is $L(L+1)/(12a)\geq L^2/(12a)$. Since
$\kappa(\sum_Ra^{-D/2})^{2/D}=L^2/a$, the geometric dependence cannot
be uniformly dropped even for uniform weights. This is not a global
counterexample to a partition algorithm or to the physical divergence-zero
observation problem.

For attribution and comparison see S.~Andres, J.-D.~Deuschel and M.~Slowik,
\emph{Harnack inequalities on weighted graphs and some applications to the
random conductance model}, arXiv:1312.5473v6, Section~2.2, (2.5)--(2.8).
Their local normalized inequalities are classical prior context; this
proof does not import their theorem, graph assumptions or Harnack conclusions.
The finite projection/entropy, coarea and power steps above are classical,
not a new general method. The gain is a complete finite derivation of the
specific input previously left external. No Lean execution, independent
review, original-priority claim, new physical hypothesis removal or new
AIRR paper follows from this supplement. The frozen PDF candidate and its
actual reports/permissions remain unchanged. This fragment is integrated
in the present changed artifact, which does not inherit their assessments.
